\chapter{Complexity — Hard Functions}
%%% generated by the blueprint pipeline from TCSlib.Complexity.CircuitComplexity.HardFunctions
\section{Overview}

Declarations of the Lean module \texttt{TCSlib.Complexity.CircuitComplexity.HardFunctions}.

\section{Declarations}

\begin{definition}[Binary encoding of a Boolean list]
\lean{BoolCircuit.bitsToNat}
\label{BoolCircuit.bitsToNat}
\leanok
For a finite list $w$ of Boolean values, map each entry to $1$ if it is true and $0$ if
it is false, append a final digit $1$, and interpret the resulting list as base-$2$
digits; the resulting natural number is the binary encoding of $w$.
\end{definition}

\begin{theorem}[Bound on Boolean digit values]
\lean{BoolCircuit.bitsToNat_digits_lt}
\label{BoolCircuit.bitsToNat_digits_lt}
\leanok
\difficulty{4}
For every finite list $w$ of Booleans, each entry of the list obtained by replacing
false by $0$ and true by $1$, then appending $1$, is less than $2$.
\end{theorem}

\begin{theorem}[Bound on the binary list encoding]
\lean{BoolCircuit.bitsToNat_lt}
\label{BoolCircuit.bitsToNat_lt}
\leanok
\uses{BoolCircuit.bitsToNat, BoolCircuit.bitsToNat_digits_lt}
\difficulty{3}
For every finite Boolean list $w$, the natural number whose base-two digits are the bits
of $w$ followed by a terminal digit $1$ is less than $2^{|w|+1}$, where $|w|$ is the
length of $w$.
\end{theorem}

\begin{theorem}[Injectivity of the binary list encoding]
\lean{BoolCircuit.bitsToNat_injective}
\label{BoolCircuit.bitsToNat_injective}
\leanok
\uses{BoolCircuit.bitsToNat, BoolCircuit.bitsToNat_digits_lt}
\difficulty{5}
The map that sends a finite Boolean list to the natural number whose base-two digits are
the list’s bits followed by a terminal digit $1$ is injective.
\end{theorem}

\begin{theorem}[Length of an encoded literal]
\lean{BoolCircuit.length_encodeCircuit_lit}
\label{BoolCircuit.length_encodeCircuit_lit}
\leanok
\uses{BoolCircuit.encodeCircuit}
\difficulty{4}
Let $n,k\in\mathbb{N}$ with $k<n$, and let $s\in\{\mathrm{false},\mathrm{true}\}$ be the
literal's polarity. The encoding of the literal with index $k$ and polarity $s$ has
length $k+3$.
\end{theorem}

\begin{theorem}[Circuit encoding length bound]
\lean{BoolCircuit.length_encodeCircuit_succ_le}
\label{BoolCircuit.length_encodeCircuit_succ_le}
\leanok
\uses{BoolCircuit.TreeCircuit, BoolCircuit.TreeCircuit.ind, BoolCircuit.TreeCircuit.size, BoolCircuit.encodeChildren, BoolCircuit.encodeChildren_cons, BoolCircuit.encodeChildren_nil, BoolCircuit.encodeCircuit, BoolCircuit.encodeCircuit_node, BoolCircuit.length_encodeCircuit_lit}
\difficulty{5}
For every Boolean circuit $C$ on $n$ variables, the length of its encoding plus one is
at most $(n+4)$ times the size of $C$.
\end{theorem}

\begin{theorem}[Injectivity of circuit encoding]
\lean{BoolCircuit.encodeCircuit_injective}
\label{BoolCircuit.encodeCircuit_injective}
\leanok
\uses{BoolCircuit.encodeCircuit, BoolCircuit.readCircuit_encodeCircuit}
\difficulty{3}
For any natural number $n$, the circuit encoding is injective: two circuit trees on $n$
variables with the same Boolean-list encoding are equal.
\end{theorem}

\begin{theorem}[Count of functions computed by small tree circuits]
\lean{BoolCircuit.card_computable_le}
\label{BoolCircuit.card_computable_le}
\leanok
\uses{BoolCircuit.TreeCircuit, BoolCircuit.TreeCircuit.eval, BoolCircuit.TreeCircuit.size, BoolCircuit.bitsToNat, BoolCircuit.bitsToNat_injective, BoolCircuit.bitsToNat_lt, BoolCircuit.encodeCircuit, BoolCircuit.encodeCircuit_injective, BoolCircuit.length_encodeCircuit_succ_le}
\difficulty{5}
For natural numbers $n$ and $S$, the number of Boolean functions on $n$ variables
computed by Boolean circuit trees of size at most $S$ is at most $2^{(n+4)S}$.
\end{theorem}

\begin{theorem}[Existence of a hard function for small circuit trees]
\lean{BoolCircuit.exists_not_eval_of_lt}
\label{BoolCircuit.exists_not_eval_of_lt}
\leanok
\uses{BoolCircuit.TreeCircuit, BoolCircuit.TreeCircuit.eval, BoolCircuit.TreeCircuit.size, BoolCircuit.card_computable_le}
\difficulty{4}
Let $n,S$ be natural numbers satisfying $(n+4)S<2^n$. Then there is a Boolean function
on $n$ input bits such that every Boolean circuit tree of size at most $S$ disagrees
with it on at least one input.
\end{theorem}

\begin{theorem}[Existence of a hard Boolean function for trees]
\lean{BoolCircuit.exists_hard_function}
\label{BoolCircuit.exists_hard_function}
\leanok
\uses{BoolCircuit.TreeCircuit, BoolCircuit.TreeCircuit.eval, BoolCircuit.TreeCircuit.size, BoolCircuit.exists_not_eval_of_lt}
\difficulty{4}
For every natural number $n$, there is a Boolean function on $n$ input bits such that
every Boolean circuit tree of size at most $2^n/(n+5)$ disagrees with the function on at
least one input.
\end{theorem}
